\centerline{\bf \Huge Домашнее задание.}
\centerline{\bf \Huge Турниры + геометрия}

\begin{flushright}
        \scriptsize{Лень чет придумывать(}
\end{flushright}

\problem{1!}
а) Сколько шахматистов играло в однокруговом турнире, если всего в этом соревновании было сыграно 190 партий?

б) Сколько шахматистов играло в однокруговом турнире, если всего было набрано больше 50, но меньше 60 очков?

в) Сколько команд играло в однокруговом турнире, если всего было набрано больше 50, но меньше 60 очков по системе 2–1–0?

\problem{2!}
Шахматист сыграл в турнире 20 партий и набрал 12,5 очков. На сколько больше он выиграл партий, чем проиграл?

\problem{3} \textbf{Лемма о перекладывании треугольников.} В треугольниках $ABC$ и $A'B'C'$, $AB = A'B'$ и $\angle BAC + \angle B'A'C' = 180^{\circ}$. Докажите, что равенство сторон $BC$ и $B'C'$ равносильно равенству углов $\angle BCA$ и $\angle B'C'A'$ .

\problem{4} В выпуклом четырехугольнике $ABCD$ известно, что $AD=BC$, $\angle ABD + \angle CDB = 180^{\circ}$. Докажите, что $\angle BAD = \angle BCD$

\problem{5} Пусть $K$ $-$ середина стороны $BC$ треугольника $ABC$. На лучах $AB$ и $AC$ взяты точки $X$ и $Y$ соответственно таким образом, что $AX = AY$ и точка $K$ лежит на отрезке $XY$. Докажите, что $BX = CY$.